\begin{homeworkProblem}[Seção 3-1]
  Seja $f:U\to\mathbb{R}^{n}$ de classe $C^{1}$ num aberto convexo $U\subseteq\mathbb{R}^{m}$, com $0\in U$ e $f(0)=0$. Se $|f^{\prime}(x)|\leq |x|$ para todo $x\in U$, então $|f(x)|\leq \frac{1}{2}|x|^{2}$ seja qual $x\in U$. Conclua que se $f(0)=0$ e $f^{\prime}(0)=0$, com $f\in C^{2}$, então $\left| \frac{\partial^{2}f }{\partial u\partial v}(x) \right|\leq |u||v|$ para $x\in U$ e $u,v\in\mathbb{R}^{m}$ quaisquer implica ainda $|f(x)|\leq\frac{1}{2}|x|^{2}$. Generalize.
  \begin{solution}
    Note que como $U$ é um aberto convexo, $f\in C^{1}$ e $f(0)=0$  pelo teorema fundamental do cálculo temos que
    \begin{align*}
      f(x)=f(x)-f(0)=\int_{0}^{1}f^{\prime}(tx)\cdot x\,dt.
    \end{align*}
    Logo temos que como $|f^{\prime}(x)|\leq |x|$ se cumpre que 
    \begin{equation}\label{eq:3-1-3}
      |f(x)|=\left| \int_{0}^{1}f^{\prime}(tx)\cdot x\,dt \right|\leq \int_{0}^{1}|f(tx)||x|\,dt\leq \int_{0}^{1}|tx||x|\,dt\leq \frac{1}{2}|x|^{2}.
    \end{equation}
    Isto, para todo $x\in U$.\\

    Agora, note que se temos que $f(0)=0$ e $f^{\prime}(0)=0$ com $f\in C^{2}$ y além mais temos que $\left| \frac{\partial^{2}f }{\partial u\partial v}(x) \right|\leq |u||v|$ para todo $x\in U$ e $u,v\in\mathbb{R}^{m}$ quaisquer, então $|f(x)|\leq\frac{1}{2}|x|^{2}$ pelo seguinte.\\
    Pelo teorema fundamental do cálculo temos que
    \begin{align*}
      f^{\prime}(x)\cdot u=(f^{\prime}(x)-f^{\prime}(0))\cdot u=\int_{0}^{1}f^{\prime\prime}(tx)\cdot (x,u)\,dt.
    \end{align*}
    Mas note que pela hipótese temos que $f^{\prime\prime}(tx)\cdot (x,u)=\frac{\partial^{2}f }{\partial x\partial u}(tx)\leq |x||u|$, assim
    \begin{equation}\label{eq:3-1-1}
      |f^{\prime}(x)\cdot u|\leq \int_{0}^{1}\left| f^{\prime\prime}(tx)\cdot (x,u) \right|\,dt\leq |x||u|.
    \end{equation}
    Logo pegando $u$ com $\left| u \right|=1$ temos que a norma do operador cumpre
    \begin{equation}\label{eq:3-1-2}
      |f^{\prime}(x)|\leq |x|.
    \end{equation}
    Logo, tamos no caso anterior y se pode conclui que $|f(x)|\leq\frac{1}{2}\left| x \right|^{2}$.\\

    \textbf{Generalização:} Seja $f:U\to\mathbb{R}^{n}$ de classe $C^{k}$ num aberto convexo $U\subseteq \mathbb{R}^{m}$, com $0\in U$ e (fazendo uso abusivo do notação) $f(0)=f^{\prime}(0)=\cdots=f^{(k-1)}(0)=0$. Se $\frac{\partial^{k}f }{\partial u_1\cdots u_k}(x)\leq |u_1|\cdots|u_k|$ para todo $x\in U$ e $u_i\in\mathbb{R}^{m}$ com $i\in\{1,\cdots,k\}$, então temos que $|f(x)|\leq \frac{1}{k!}|x|^{k}$.\\
      Então, note que fazendo o mesmo cálculo que \eqref{eq:3-1-1} e \eqref{eq:3-1-2} temos que
      \begin{align*}
        |f^{(k-1)}(x)|\leq |x|.
      \end{align*}
      Logo faendo o mesmo cálculo que \eqref{eq:3-1-3}, temos que isto implica que
      \begin{align*}
        |f^{(k-2)}(x)|\leq \frac{1}{2}|x|^{2}.
      \end{align*}
      Logo fazendo um raciocinio indutívo temos que
      \begin{align*}
        |f(x)|\leq\frac{1}{k!}|x|^{k}.
      \end{align*}
      (Intuição: Note que para cada derivada estamos ganhando um $|x|$, portanto estamos ganhando uma potência em cada derivada faltante; além disso, por essa mesma potência, temos que a integral de \eqref{eq:3-1-3} o valor de $t$ ganha essa mesma potência, de modo que cada integral resulta em uma constante $\frac{1}{j}$, portanto a multiplicação dessas constantes é o fatorial das derivadas faltantes).
  \end{solution}
\end{homeworkProblem}
